\documentclass[11pt]{article}
\usepackage[a4paper,margin=28mm]{geometry}
\usepackage{amsmath,amssymb,amsthm}
\usepackage[hidelinks]{hyperref}
\newtheorem{proposition}{Proposition}
\newcommand{\Aut}{\operatorname{Aut}}
\title{Chirality without firmness}
\author{Human Compression project}
\date{}
\begin{document}
\maketitle

In the final published edition of Leemans--Tranchida \cite{LT}, printed
page 492, a geometry is an incidence system in which every flag extends
to a chamber. Firmness is defined separately. On page 493, chirality
requires exactly two orbits of the full type-preserving group on
chambers, with adjacent chambers in different orbits. These definitions
permit singleton panels. The following example shows why the omission
matters for the unrestricted auxiliary statements.

\section{An example with seven objects}
Let the three type classes be
\[
 X_0=\{d,c_1,c_2\},\qquad X_1=\{x_1,x_2\},\qquad
 X_2=\{y_1,y_2\}.
\]
Apart from reflexivity, declare precisely the following pairs incident,
for $j=1,2$:
\[
 x_j\sim y_j,\qquad d\sim x_j,\quad d\sim y_j,
 \quad c_j\sim x_j,\quad c_j\sim y_j.                 \tag{1}
\]
Call the resulting system $\Gamma$.

\begin{proposition}
The system $\Gamma$ is a connected geometry satisfying the literal
chirality definition of \cite{LT}. Its full type-preserving group is
$C_2$. It is not firm. Moreover, all residues of corank at least two
at the chamber $C_1=\{c_1,x_1,y_1\}$ are connected, but the residue
of $d$ is disconnected.
\end{proposition}
\begin{proof}
The only chambers are
\[
 C_j=\{c_j,x_j,y_j\},\qquad D_j=\{d,x_j,y_j\}
 \quad (j=1,2).
\]
Every pair in (1) lies in one of these chambers, and every object
lies in a pair. Hence every flag extends to a chamber. The whole
incidence graph is connected: $d$ is joined to every $x_j,y_j$, and
$c_j$ is joined to $x_j$.

Among the type-zero objects, $d$ has degree four and the other two
have degree two. Every type-preserving automorphism fixes $d$.
The matching $x_j\sim y_j$ forces the same permutation of the two
indices on types one and two. The incidences of $c_j$ then force
that same permutation on $c_1,c_2$. Conversely, simultaneous exchange
of the two indices preserves all incidences. Thus the full group is
exactly $C_2$, with chamber orbits $\{C_1,C_2\}$ and $\{D_1,D_2\}$.
The only adjacent chamber pairs are $C_j,D_j$, so adjacency always
crosses between the two orbits. This is the stated chirality condition.

The panel $\{c_j,x_j\}$ has only the completion $y_j$, so the
geometry is not firm. In fact eight panels have one completion and
two panels have two completions.

The flags in $C_1$ of corank at least two are the empty flag and
the three singleton flags. The empty residue is the connected whole
geometry. The residue of $c_1$ is the single edge $x_1\sim y_1$.
The residue of $x_1$ has vertices $c_1,d,y_1$ and edges
$c_1\sim y_1\sim d$; the residue of $y_1$ is the analogous path.
All four residues are connected. In contrast, the residue of $d$
has vertices $x_1,x_2,y_1,y_2$ and consists of the two disjoint edges
$x_1\sim y_1$ and $x_2\sim y_2$.
\end{proof}

The example refutes three literal auxiliary assertions: the inference
from chirality to thinness on printed page 493; Theorem 4.1 on page
494, applied to $C_1$; and Proposition 4.2 on page 495. For the last
assertion, already the proper flag type $\{0\}$ has two object orbits,
$\{d\}$ and $\{c_1,c_2\}$. The source proofs use an adjacent chamber
where firmness has not been supplied. With firmness as an additional
hypothesis, both normalization arguments are valid and are retained
in the project proof bank.

This example is not a standard $C^+$ coset geometry. Indeed no subgroup
of its full type-preserving group is transitive on the type-zero
objects, whereas a coset action is transitive on each type. Thus it
does not refute the corrected higher-rank variant of Corollary 4.9.
The separate cyclic example handles the unrestricted wording of that
corollary in rank two.

\section{The rank of the residue in Proposition 4.7}
The displayed statement of Proposition 4.7, printed page 496, requires
$|J|<|I|$ and asserts connectedness of the residue under
\[
 G_J=\langle G_{J\cup\{i\}}:i\in I\setminus J\rangle.
                                                               \tag{2}
\]
The preceding discussion excludes rank-one residues as irrelevant to
residual connectedness. This contextual restriction must also be
present in the proposition: (2) does not imply connectedness in rank
one. For example, in the already accepted rank-four $C^+$ geometry
of \cite{HC}, take any base panel with type set $J$, so $|J|=3$.
The intersection property gives $G_J=G_I=1$, and hence (2) holds.
But the residue consists of two objects of the same type. Its
incidence graph has two isolated vertices and is disconnected.
The existing recovery of this sufficient criterion is explicitly for
residues of rank at least two and for the designated action with
proper-flag transitivity. The source's separate switch between the
full automorphism group and the designated group is documented in
the accepted principal counterexample; neither qualification is
silently removed here.

\section{Verification and scope}
The accompanying finite check enumerates every flag of the seven-object
system and all $3!\,2!\,2!=24$ permutations preserving its types.
It finds exactly the two automorphisms proved above, the two chamber
orbits, the complete panel counts and the four connected base residues.
The disconnected residue of $d$ is checked separately. No group search,
classification or automorphism software is required.

The definitions and statements belong to \cite{LT}; the seven-object
construction and proof are supplied here without a priority claim.
The rank-one observation imports the reviewed data of \cite{HC}.
These auxiliary corrections add no original, open-question solution,
Lean theorem or compression measurement. The completion ledger keeps
the literal assertions separate from their corrected variants.

\begin{thebibliography}{2}
\bibitem{LT} D.~Leemans and P.~Tranchida,
\emph{On residual connectedness in chiral geometries},
Algebraic Combinatorics \textbf{4} (2021), 491--499,
\href{https://doi.org/10.5802/alco.162}{doi:10.5802/alco.162}.
\bibitem{HC} Human Compression project,
\emph{A chiral coset geometry with disconnected residues},
reviewed project manuscript, five pages; exact proof and review bound
in the dependency ledger.
\end{thebibliography}
\end{document}
